This chapter synthesizes prior research relevant to user–environment coupling in and around head-mounted wearables (HMWs) and earables. We organize the literature into four areas: (1) Speech and Conversation Enhancement, (2) Human Activity Recognition, (3) Acoustic Sensing Systems, and (4) Other Novel Applications on Head-mounted Wearables. Each area reviews representative methods, their applicability to HMWs, and open gaps that motivate our contributions.

\section{Speech and Conversation Enhancement}

\subsection{Speech Enhancement}
Speech enhancement is a longstanding problem in voice communication and assistive listening. Methods can be grouped by their primary approach: \textit{model-driven}, \textit{machine learning}, \textit{target/user-centric}, and \textit{multi-modal}.

\subsubsection{Model-driven}
Classical approaches assume stationarity/independence in the time–frequency domain to separate speech from noise \cite{ephraim1995signal}, or leverage multi-microphone arrays and beamforming to steer toward a talker for enhancement \cite{yousefian2014hybrid,ji2017coherence} and separation \cite{wood2017blind,kitamura2016determined}. While effective under controlled assumptions, these methods struggle with dynamic, non-stationary noise and with the compact, irregular microphone layouts typical of HMWs.

\subsubsection{Machine learning}
Deep neural networks (DNNs) have advanced speech enhancement for both single-channel \cite{subakan2021attention} and multi-channel audio with arbitrary layouts \cite{zhang2021microphone}. For instance, ClearBuds \cite{chatterjee2022clearbuds} processes stereo earbud audio with a DNN, but introduces communication overhead and degrades when competing sources share the same angle. Despite strong results, audio-only DNNs depend on training-domain coverage and can underperform with slim-volume, distant-mic recordings common in HMWs.

\subsubsection{Target/user-centric enhancement}
Beyond location-based steering, user-centric enhancement aims to selectively extract desired sounds while suppressing others, enabling superhuman listening when end-to-end latency is within \mbox{20–30 ms} \cite{stone1999tolerable, gupta2020acoustic}. Fixed-target suppression (e.g., classical noise reduction \cite{wiener1949extrapolation}) assumes stationary noise, whereas dynamic target extraction can condition on speaker embeddings \cite{zmolikova2023neural, veluri2024look}, articulatory motion (e.g., lip vibration) \cite{liu2021wavoice}, or spatial constraints (e.g., SoundBubble) \cite{chen2024hearable}. Other systems attenuate non-speech by class \cite{veluri2023semantic}, via learned audio embeddings \cite{liu22w_interspeech}, or by reusing prior ambient audio for cancellation \cite{shen2018mute}. These designs typically emphasize single-user scenarios or fixed categories, leaving gaps for dynamic, multi-party conversations where both verbal and non-verbal cues shift rapidly.

\subsubsection{Multi-modal}
Multi-modal designs incorporate visual, wireless, or inertial signals to stabilize enhancement.

\paragraph{Audio–Visual}
Camera streams can guide separation and enhancement by correlating lip motion and articulation with audio \cite{michelsanti2021overview, gao2021visualvoice}. These approaches leverage deep cross-modal embeddings to extract target speech from noise but require a visual sensor, which is not always available or practical on wearables.

\paragraph{Audio–Wireless}
Ultrasonic probing transmits inaudible signals (e.g., >17 kHz) and analyzes echoes from articulators to enhance noisy audio \cite{sun2021ultrase, zhang2021sensing}. Coupled mmWave and microphones can aid speech-related tasks \cite{liu2021wavoice, ozturk2021radiomic}, but mmWave hardware is bulky and uncommon on HMWs, and ultrasonic methods often constrain user posture (e.g., a phone fixed toward the mouth).

\paragraph{Audio–IMU}
Bone conduction and accelerometers provide user-centric vibration signals that correlate with the wearer’s voice. Fusing these with air microphones improves robustness \cite{wang2022fusing, tagliasacchi2020seanet}, yet bandwidth limits on commodity sensors and anatomical variability challenge generalization \cite{Inertial63:online}. System evaluations also need real-world noises and unseen users to validate practicality.

Non-acoustic sensing (IMUs, piezo) can capture contact sounds \cite{maruri2018v, khanna2021jawsense}; remote vibrometry via geophones/LiDAR/mmWave detects limited speech content \cite{han2017pitchln, sami2020spying, li2020vocalprint}, but cannot reconstruct full-spectrum audio \cite{anand2018speechless}. Location-based audio augmentation (e.g., museum guides \cite{harma2004augmented}) and Ear-AR \cite{yang2020ear} demonstrate spatial alignment and head-tracked rendering, yet assume predefined triggers and static contexts, not highly dynamic group conversations.

\subsection{Conversation Modeling}
Natural conversation blends verbal content with non-verbal behavior. For online enhancement and social awareness, both streams matter.

\paragraph{Non-verbal feature}
Non-verbal features, or body language \cite{d2010multimodal}, play a key role in conversations. Ego-centric tasks such as “look at me” and “talk to me” \cite{grauman2022ego4d} benefit from multi-modal glasses \cite{donley2021easycom}, yet remain challenging with audio alone. Conventional localization emphasizes direction of arrival \cite{wang2022nerc, yang2022deepear}, which misses speaker orientation—a crucial cue for “talking to me.” Acceleration-based detection \cite{gedik2018detecting} captures the wearer’s motion but lacks broader context. Joint estimation of location and orientation \cite{ahuja2020direction} typically assumes static arrays, limiting adaptability in dynamic HMW settings.

\paragraph{Verbal feature}
Speaker diarization \cite{bredin2023pyannote} and conversational analysis \cite{wei2022conversational} support tasks like meeting summaries \cite{ramprasad2024analyzing} and readability improvements \cite{wang2024diarizationlm}. Turn-taking shows clear temporal patterns \cite{levinson2015timing} and has been used for proactive mediation \cite{lee2018flower} and event detection via overlaps/backchannels \cite{chen2024target, chang2022turn}. Yet turn-taking cues are only reliable after events occur, introducing latency for online enhancement. Complementary approaches—including LLM-guided methods (e.g., SocialMind \cite{yang2025socialmind})—combine non-verbal and verbal cues. Overall, audio-only modeling needs richer non-verbal proxies and low-latency signals tailored to mobile HMWs.

\paragraph{Conversation Enhancement}
Conversation enhancement aims to amplify the speech that matters to the wearer while preserving social context in dynamic, multi-party settings. It must decide who to enhance, whether that speaker is addressing the wearer, and where they are located, then render or suppress sources accordingly within the ANC latency budget (\mbox{20–30 ms}) \cite{stone1999tolerable, gupta2020acoustic}. Practical pipelines combine low-latency non-verbal cues (e.g., head orientation, proximity, motion) with verbal cues (e.g., diarization and speaker embeddings \cite{bredin2023pyannote, zmolikova2023neural, veluri2024look}) to gate, separate, and spatialize streams. User-centric signals such as bone-conducted vibration, in-ear occlusion, and facial acoustics \cite{he2023towards, han2024earspeech, zhang2025wearse} can condition target extraction on the wearer’s own voice, while audio-only enrollment can provide embeddings for companions. Multi-device collaboration (WASN) augments spatial context when available \cite{gburrek2021geometry}, but must balance synchronization, privacy, and compute. Prior systems like SoundBubble \cite{chen2024hearable} constrain geometry. A conversation-driven design should adapt online to turn-taking and orientation, avoid camera dependencies, and maintain energy/compute budgets on HMWs.

\section{Human Activity Recognition}

\subsection{Motion sensing with IMUs}
IMUs in mobile and wearable devices support activity and motion recognition, with performance sensitive to placement. Examples range from smartphone-based recognition (e.g., LIMU-BERT \cite{xu2021limu}) to smartwatch sensing \cite{chan2024capture}. Large datasets still yield coarse-grained labels (activity level, daily activities) \cite{chan2024capture}. Earables offer more stable placement \cite{lyu2024earda, han2023headmon, han2023headsense, han2025headmon+}, but face similar limits in fine-grained HAR.
Multiple-IMU deployments can recover body pose (e.g., up to 17 IMUs) and enable fine-grained recognition (MotionBERT \cite{zhu2023motionbert}), yet are impractical at scale. Combining commodity devices (phone/watch/earables) improves granularity \cite{stromback2020mm} and can estimate pose with three IMUs (IMUPoser \cite{mollyn2023imuposer}), but at higher cost and setup complexity. Pretrained models linking IMU to text/vision enable zero-shot recognition (IMU2CLIP \cite{moon2023imu2clip}, ImageBind \cite{girdhar2023imagebind}), and recent LLM-based methods reason over motion data \cite{ji2024hargpt, leng2024imugpt, xu2024autolife, ouyang2024llmsense, xu2024penetrative}, though fine-grained accuracy remains a challenge due to limited environmental/context cues.

\subsection{Acoustic sensing for HAR}
Passive audio supports speech recognition and sound classification \cite{schmid2023efficient}, but maps only indirectly to human activity and is confounded by loudspeakers \cite{cristina2024audio}. Multi-channel passive sensing estimates direction of arrival with arrays \cite{schmidt1986multiple, adavanne2018sound} or binaural earables \cite{yang2022deepear, krause2021joint}, improving spatial context but not confirming human presence. To detect human vocal activity, systems use feedforward microphones \cite{zhang2024earsavas} or IMU bone conduction \cite{he2023towards}, which are specific to vocal events. Active acoustic sensing (probing and echo analysis) expands capability: facial motion \cite{dong2024rehearsse}, upper-body pose \cite{mahmud2023posesonic}, and gestures \cite{yang2024maf}. However, speaker placement and privacy constraints of wearables limit outward acoustic probing range and fidelity.

\subsection{Multi-modal sensing for HAR}
Audio–IMU fusion addresses modality gaps and boosts robustness. Earable systems combine bone-conducted vibration with air microphones to separate the wearer’s speech from others \cite{he2023towards}. Other multi-device wearables combine IMU and audio for richer activity sets \cite{min2018exploring, mollyn2022samosa, bhattacharya2022leveraging}. Yet multi-modal fusion is not always superior: on diverse smart-glasses datasets, multi-modal models can underperform single-modality specialists \cite{gong2023mmg, huang2022modality}. Incorporating external knowledge (context triggers, LLMs, cross-modal embeddings) helps \cite{mollyn2022samosa, yang2025contextagent, xu2024autolife, han2023onellm, girdhar2023imagebind, yang2024drhouse}, but an explicit, effective fusion of heterogeneous mobile data remains open.

\section{Acoustic Sensing Systems}

\subsection{Audible and ultrasonic}
Audible sensing derives semantics (transcription \cite{gao2023funasr}, classification \cite{gemmeke2017audio}) and spatial cues from multi-channel recordings \cite{wang2019doa}. Ultrasonic sensing transmits inaudible probes to communicate and sense, supporting underwater localization/communication \cite{chen2022underwater, chen2023underwater}, backscatter \cite{jang2019underwater}, air–water communication \cite{tonolini2018networking}, SOS beacons \cite{yang2023aquahelper}, dual-channel communication \cite{qian2021aircode}, gesture \cite{wang2020push}, vital signs \cite{wang2022loear}, and metasurface-enhanced interactions \cite{zhang2023acoustic}. While ultrasonic probes unlock controllable measurements, they require hardware access and careful design to satisfy safety, privacy, and energy constraints on HMWs.

\subsection{Wireless acoustic sensor networks (WASN)}
Moving beyond single-device processing, WASN treats distributed microphones as a cooperative network. With ambient sound, nodes can be geometrically calibrated \cite{gburrek2021geometry} using DoA \cite{wang2019doa}, distance \cite{gburrek2021iterative}, or both \cite{gburrek2021geometry}, solved by iterative optimization \cite{fischler1981random} or dataset matching \cite{hennecke2009hierarchical}. Additional challenges include sample rate offset (SRO) \cite{gburrek2022synchronization}, distributed beamforming \cite{gburrek2023integration}, and swarm-style acoustics \cite{itani2023creating}. For HMWs, ad-hoc, on-body/off-body collaboration promises richer spatial reasoning but must balance synchronization, privacy, and compute.

